% !TEX program = lualatex
\documentclass[aspectratio=43,professionalfonts,handout]{beamer}
\usepackage{beamer_template}

\newcommand{\LectureNum}{9}
\newcommand{\LectureTitle}{一般の周期関数のフーリエ級数（1）}
\newcommand{\TextPages}{p.\,80-85}
\newcommand{\CourseName}{応用数学}
\renewcommand{\TermName}{前期}

\begin{document}

\begin{frame}{\CourseName\quad 第\LectureNum 回「\LectureTitle」}
  {\small \TermName\quad \TeacherName\quad ／\quad 教科書 \TextPages}

  \textbf{目標}：周期2lの周期関数のフーリエ級数を求めることができる

\end{frame}

\begin{frame}{周期2lへの一般化}
  \begin{block}{}
  \begin{itemize}
    \item 変数変換 $t = \pi x/l$
    \item フーリエ係数の一般形
    \item $l=\pi$ のとき第 8 週の公式に一致
  \end{itemize}
  \end{block}
\end{frame}

\begin{frame}{例題2の計算}
  \begin{block}{}
  \begin{itemize}
    \item 周期2の矩形波
    \item $c_{0}, a_{n}, b_{n}$ を計算
    \item フーリエ級数の表示とグラフ
  \end{itemize}
  \end{block}
\end{frame}

\begin{frame}{偶関数・奇関数の性質}
  \begin{block}{}
  \begin{itemize}
    \item 偶関数：$b_{n}=0 \to $ フーリエ余弦級数
    \item 奇関数：$c_{0}=a_{n}=0 \to $ フーリエ正弦級数
    \item 計算量が半分になる
  \end{itemize}
  \end{block}
\end{frame}

\begin{frame}{例題3・4}
  \begin{block}{}
  \begin{itemize}
    \item f(x)=|x|：偶関数  $\to$  余弦級数
    \item f(x)=x(1-x)：奇関数  $\to$  正弦級数
  \end{itemize}
  \end{block}
\end{frame}

\begin{frame}{演習}
  \begin{exampleblock}{自分で解いてみよう}
  \begin{itemize}
    \item 問2〜4より選題
  \end{itemize}
  \end{exampleblock}
\end{frame}

\begin{frame}{まとめ}
  \begin{itemize}
    \item 周期2lの公式
    \item 偶関数・奇関数の活用
    \item 次回：一般の周期関数（2）・収束定理
  \end{itemize}
\end{frame}

\end{document}
